\section{Discussion}
\label{sec:disc}
The significant performance benefits of Vertex-Fed are driven by two main factors. First, the gradient cosine similarity check allows the network to isolate nodes that possess heavily corrupted or adversarial data points. During the training run, we observed that malicious nodes or clients experiencing high measurement noise were naturally isolated from the main topology by the Nimbus controller.

Second, the dynamic pruning scheme reduces consensus time on the graph. A well-known limitation of decentralized optimization is the consensus delay term, which depends on the second-largest eigenvalue of the consensus matrix. By dynamically optimizing the topology to prune slow communication paths while retaining high-capacity channels, Vertex-Fed improves the spectral properties of the consensus graph, accelerating convergence.

\subsection{Limitations}
A potential limitation of our framework is the need to periodically re-evaluate pairwise similarity metrics, which introduces an $\mathcal{O}(M^2)$ computational complexity at the controller. However, in practice, the similarity evaluation interval $\tau$ can be set to relatively large values (e.g., every 50 epochs) without hurting convergence stability, which limits the computational overhead.
